\chapter{Discussion}

Reviewing the trends in the models, we recap several important trends about the methodologies used and their results. The linear regressions and regularized regressions achieved moderate $R^2$ and Adjusted-$R^2$ values for predicting U.S. returns using only U.S. economic data and for predicting Chinese returns using only Chinese economic data, achieving $0.70\leq R^2\leq 0.80$ for the U.S. market and $0.55\leq R^2\leq 0.60$ for the Chinese market. While the Chinese economic data gave no additional explanatory power to the U.S. market models, the U.S. economic data more consistently improved the predictions of Chinese returns. Still, there were some issues with the linear models, with residual analyses revealing heteroskedasticity and otherwise suggesting nonlinearity in the factors.

The SpAM models often saw similar $R^2$ and Adjusted-$R^2$ values, showing that the generalization beyond linearity did not hurt the models ability to pick up on trends. In some cases, the $R^2$ and Adjusted-$R^2$ values were even higher; however, the model was more prone to overfitting on the data, especially in cases where the time frames were short or when using larger number of predictors. The model suffered from the lack of interpretability as well, where a plot was required to understand the relationship between an economic measure and the returns.

The kernel regression models often had very high $R^2$ and Adjusted-$R^2$ scores, indicating that the model could almost replicate the in-sample data. This was, in a lot of cases, the result of severe overfitting. Still, there were some interesting results captures in some of the economic curves, and even more interesting results in cross-country prediction. Though consistent with previous approaches, cross-country prediction trends were very interesting to examine in the nonparametric model. The partial dependence plots revealed some interesting trends in the factors that aligned with economic intuition and were very helpful in assessing overfitting. 

As far as general trends within the results, domestic economic data was the strongest predictor for both the U.S. and Chinese market returns. The Chinese market gained more explanatory power when U.S. economic data was introduced, which supports the view that global markets receive stronger signals from the U.S. market than in the opposite direction.

The models with the least assumptions also had a strong tendency to overfit to the data, especially in scenarios where there were limited samples and a larger number of predictors. While logical in principle, this posed a great challenge when analyzing market data, especially when segmenting for structural breaks in the data.

Throughout the models, there is a trade off between flexibility and the ease of interpretation. Linear models worked really well, achieving high $R^2$ values while maintaining a good level of easy interpretability. As the models became less structures, and thus more capable of handling more complex data, they also became more difficult to interpret. The partial effect plots in the SpAMs and kernel regressions also suffered from interpretability as it makes it even harder to understand how a factor is impacted with other factor movement.

When using less-structured models, it is more difficult to get enough data without breaking the stationarity assumption in the data. Assessments of structural breaks in the market yielded several scenarios over a few decades where the underlying structural factors had shifted behind the data. This is a significant challenge when attempting to understand the intricacies of the market.

There are several ways to build off this work. The first next step is to introduce models that can naturally capture interactions among the economic factors without requiring predefined functional forms, perhaps with tree-based methods. Alternatively, using the several decades of data available to better estimate the GAMs and SpAMs to allow for pairwise interactions. In order to deal with structural breaks in the data, one can try to implement indicator variables in their analyses. If attempting to work on SpAMs or kernel regressions, more nuanced regularization can help mitigate the large amount of overfitting seen in the models.

Another natural first step would be to attempt the same methods on rolling-window estimates, which might capture the trends of the data better and reduce the noise in the model estimates. Analyses in individual sectors might benefit from stronger and more interpretable insights as sector-specific spillovers are easier to detect, model, and understand.

The economic signals picked up can also be tested for their market power, creating specific portfolio strategies using the indicators from each country. This would be an interesting way to quantity the strength of certain signals and the benefits of accounting for information from other markets.

The relationship between the U.S. and Chinese markets has been long studied by researchers. The empirical challenges of quantifying such a relationship have changed significantly over the past decade. Economic and financial data are now publicly accessible at a rate greater than ever before, allowing for exploration in the markets.

The research done in this thesis should serve as a stepping stone for testing for market integrations and cross-market signals. Mis-specified models and overfit partial effects provide more information about the models and about their relations to the markets. With global markets trending towards increasing co-integration, they will reveal more and more information about their intricacies as time goes by. 